\paragraph{Harness-level retained objects.}
The preceding mechanisms can also be viewed as edits to an executable harness rather than as isolated memories or policies. Context-engineering systems update the instructions, evidence structures, and context-management procedures presented to the model~\cite{zhang2026ace,ye2026mce}; workflow-search systems modify the code-level arrangement of model calls, tools, branches, and verification steps~\cite{hu2025adas,zhang2025aflow}; and broader self-harness systems propose persistent changes to prompts, tools, middleware, skills, sub-agent configurations, or long-term memory from execution traces~\cite{zhang2026selfharness,lin2026ahe}. The persistent object is then a versioned \emph{harness patch} together with the failure evidence, intended effect, and regression risk that motivated it.

This broader view does not enlarge the survey to all agent engineering. A harness update belongs to Feedback Learning here only when it changes future retrieval-grounded multi-hop behavior, such as query progression, source selection, context preservation, verification, or stopping. A coding-agent edit that improves repository manipulation without an information-acquisition claim remains outside the core method taxonomy. For included work, the required evidence is stricter than a higher post-edit score: the edited component must be identified, its invocation or use must be observable, and gains must survive held-out regression tests under fixed model, tool, evaluator, and budget conditions.
